% Workshop project - ไฟล์หลักของโครงงาน (โครงที่คอมไพล์ได้ แต่ยังไม่มีเนื้อหา)
% เปลี่ยนชื่อเรื่อง เพิ่มเนื้อหาในโฟลเดอร์ sections/ และเพิ่มรายการใน refs.bib
% เอกสารภาษาอังกฤษ   latexmk -pdf main.tex
% เอกสารภาษาไทย      ลบ % หน้าบรรทัด thaispec ด้านล่าง แล้วสั่ง  latexmk -xelatex main.tex
%                    (ต้องมีไฟล์ Laksaman.fontspec อยู่ในโฟลเดอร์เดียวกัน)
\documentclass[11pt]{article}

% ---------------------------------------------------------------- packages
% \usepackage[thaifont=Laksaman]{thaispec}
\usepackage{amsmath}
\usepackage{amssymb}
\usepackage{amsthm}
\usepackage{graphicx}
\usepackage[backend=biber, style=numeric]{biblatex}
\addbibresource{refs.bib}
\DeclareLanguageMapping{thai}{english}   % Thai documents: bibliography words (In, pp.) stay English
\graphicspath{{figures/}}

% ---------------------------------------------------------------- theorems
\theoremstyle{plain}
\newtheorem{theorem}{Theorem}[section]
\theoremstyle{definition}
\newtheorem{definition}[theorem]{Definition}
\newtheorem{example}[theorem]{Example}

% ---------------------------------------------------------------- macros
% สร้างคำสั่งของตนเองอย่างน้อย 2 คำสั่ง แล้วใช้ในเนื้อหา
\newcommand{\R}{\mathbb{R}}

\title{Project Title}
\author{Author Name}
\date{February 2027}

\begin{document}

\maketitle

\begin{abstract}
Write a short summary of the whole document here.
\end{abstract}

\tableofcontents

% Introduction: why the topic matters and how the document is organised.
\section{Introduction}\label{sec:introduction}
Replace this paragraph. Cite sources from refs.bib like this
\cite{lamport1994latex}, and refer to later sections like this:
Section~\ref{sec:body}.

% The main part. Split it into more files (and more \input lines in main.tex) if you like.
% It must contain a figure and a table (with captions and labels), numbered and
% multi-line equations, and at least one definition, theorem or example.
\section{Main Part}\label{sec:body}
Replace this section with your content.

% Conclusion: what the reader should remember.
\section{Conclusion}\label{sec:conclusion}
Replace this paragraph.


\printbibliography

\section*{AI Use Statement}
State which AI tools you used, for what, and what you checked yourself.
If you used none, write: No AI tools were used to write this document.

\end{document}
